\subsection{Федеральный закон № 149-ФЗ «Об информации, информационных технологиях и о защите информации»}

Федеральный закон № 149-ФЗ «Об информации, информационных технологиях и о защите информации» является базовым нормативным правовым актом, регулирующим отношения, возникающие при осуществлении права на поиск, получение, передачу, производство и распространение информации, а также при применении информационных технологий и обеспечении защиты информации. Данный закон конкретизирует конституционные положения в сфере обращения с информацией и определяет правовые основы её использования в информационных системах.

В соответствии со статьей 2 Федерального закона № 149-ФЗ под информацией понимаются сведения (сообщения, данные) независимо от формы их представления, а информационная система рассматривается как совокупность содержащейся в базах данных информации и обеспечивающих её обработку информационных технологий и технических средств. Указанные определения имеют непосредственное значение для разрабатываемого программного средства, поскольку Kubernetes-кластер и связанные с ним audit-логи могут рассматриваться как часть информационной системы, в которой осуществляется сбор, хранение и обработка данных о действиях пользователей и сервисных учётных записей.

Статья 5 указанного закона устанавливает, что информация может являться объектом публичных, гражданских и иных правовых отношений, а также может свободно использоваться и передаваться, если федеральными законами не установлены ограничения доступа к ней либо иные требования к порядку её предоставления или распространения. Для рассматриваемой работы это означает, что audit-логи могут использоваться в целях мониторинга и анализа только при соблюдении установленного режима доступа и требований к защите сведений, содержащих ограниченно доступную информацию.

Особое значение для темы исследования имеет статья 16 Федерального закона № 149-ФЗ, согласно которой защита информации представляет собой принятие правовых, организационных и технических мер, направленных на обеспечение защиты информации от неправомерного доступа, уничтожения, модифицирования, блокирования, копирования, предоставления и распространения, а также на соблюдение конфиденциальности информации ограниченного доступа. Следовательно, при разработке системы обнаружения аномалий в Kubernetes-кластере необходимо предусмотреть меры по разграничению доступа к логам, защите каналов передачи данных, контролю целостности хранимой информации и ограничению несанкционированного использования результатов анализа.

Таким образом, Федеральный закон № 149-ФЗ формирует правовую основу для обработки audit-логов Kubernetes-кластера как информации, подлежащей правовому режиму и защите. При проектировании и внедрении разрабатываемого решения должны быть обеспечены законность обработки данных, соблюдение ограничений доступа, а также применение организационных и технических мер защиты информации в соответствии с действующим законодательством Российской Федерации.